% Inactive research fragment. Not an author input or a novelty claim.
% The analytic arguments use smooth potentials until capacity is established.
\section{Smooth local capacities and decreasing Hessian products}

Let $U$ be a coordinate ball with compact closure in a holomorphic
chart, and let $\omega$ be the restriction of the Hermitian metric.
Throughout, $2\leq m<n$, as in the Hessian chapter. A smooth local function is
$m$-subharmonic relative to $\omega$ when
$\ddc f\in\overline{\Gamma_m(\omega)}$ pointwise. Put
\[
 \Lambda_s(f)=(\ddc f)^s\wedge\omega^{n-s},\qquad 0\leq s\leq m,
\]
where $\Lambda_0(f)=\omega^n$. All products before the rough-coefficient lemma are
ordinary smooth products. No rough local Hessian construction or
singular weighted Sobolev space is assumed.

\begin{lemma}\label{draft621:local-smooth-bounds}
Fix $K\Subset U$ and $M<\infty$. For smooth local $m$-subharmonic
$f_1,\ldots,f_q$ with $|f_i|\leq M$, $0\leq q\leq m$, their mixed
trace has mass at most a constant depending only on $K,U,\omega,m,M$.
If also $f$ is smooth in this class and $|f|\leq M$, then, for
$q\leq m-1$,
\[
 \int_K\mathrm df\wedge\dc f\wedge
       \ddc f_1\wedge\cdots\wedge\ddc f_q
                         \wedge\omega^{n-q-1}\leq C.
\]
\end{lemma}
\begin{proof}
For one potential and a nonnegative compactly supported cutoff $\chi$,
closedness of $\ddc f$ gives
\[
 \int\chi(\ddc f)^q\omega^{n-q}
       =\int f(\ddc f)^{q-1}\ddc(\chi\omega^{n-q}).
\]
For $q\geq1$, the absolute-form estimate
\cref{hessian:absolute-mixed-forms} applies at degree $q-1<m$.
Nested cutoffs and induction, starting with volume, therefore bound
the left side by $C M^q$; replacing $M$ by $1+M$ includes $M=0$.
For mixed factors, positivity of every term in the expansion of
$(\sum_i\ddc f_i)^q\omega^{n-q}$ bounds the mixed trace by a fixed
multiple of the single trace for $\sum_i f_i$. This proves the first
assertion, including its degree-zero convention.

Replace $f$ by $f+M$, so that $0\leq f\leq2M$ without changing its
gradient or Hessian. The identity
\[
 \ddc(f^2)=2f\ddc f+2\mathrm df\wedge\dc f
\]
also shows that $f^2$ belongs to the local cone. Wedge with the
remaining $q$ Hessians and $\omega^{n-q-1}$. The first term on the
right is nonnegative, and the mixed mass estimate applied to $f^2$
and the other potentials bounds the second term. The degree used is
$q+1\leq m$. A cutoff equal to one on $K$ gives the stated local
estimate.
\end{proof}

For Borel $E\subset U$ define the smooth capacity
\[
 c_s(E)=\sup\left\{\int_E \Lambda_s(a):
        a\in C^\infty(U),\ -1\leq a\leq0,
        \ddc a\in\overline{\Gamma_m(\omega)}\right\},
       \qquad 0\leq s\leq m-1.
\]
In particular $c_0(E)=\int_E\omega^n$. Each $c_s$ is monotone,
countably subadditive and continuous under increasing unions.
For the last assertion, the supremum over competitors commutes with
the increasing supremum over sets. Local smooth mass bounds make it
finite on compact subsets. Its outer-open version is
\[
 c_s^*(E)=\inf\{c_s(G):E\subset G\subset U,\ G\text{ open}\}.
\]
It is monotone and countably subadditive and agrees with $c_s$ on
open sets. A function on a compact $K\Subset U$ is quasicontinuous
for this capacity if, for every $\epsilon>0$, it is continuous on
$K\setminus G$ for some open $G\subset U$ with $c_s(G)<\epsilon$.

The capacities control smooth mixed traces uniformly in the
potentials. Choose a fixed smooth strictly plurisubharmonic $b$ with
$\ddc b\geq\omega$ on $U$. For $q\leq s$, translate and divide
$f_1+\cdots+f_q+b$ by a fixed $L\geq1$ so that it lies in $[-1,0]$.
The constant $L$ depends only on the common bound and the fixed data.
Positivity of the smooth binomial expansion gives
\begin{equation}\label{eq:draft621:smooth-capacity-bound}
 (\ddc f_1+\cdots+\ddc f_q)^q\omega^{n-q}
       \leq C \Lambda_s\bigl((\textstyle\sum_i f_i+b-B)/L\bigr).
\end{equation}
For $q=0$ the left side is $\omega^n$; the same argument applies.
Indeed, replacing any of the $s-q$ factors $\ddc b$ by $\omega$
subtracts a nonnegative mixed trace. Thus every such mixed trace,
and every coefficient of a smooth product of $q<m$ Hessians, is
dominated on Borel sets by $C c_s$. The latter assertion follows
from the absolute mixed-form estimate.

\begin{lemma}\label{draft621:smooth-torsion}
Let $0\leq q\leq m-2$ and let $S$ be a smooth real $(1,1)$-form
in $\overline{\Gamma_m(\omega)}$ pointwise.
Then $T=S^q\omega^{n-q-2}$ is strongly positive. For smooth real $f$
and nonnegative Borel weights $a$, supported on a fixed compact set,
\[
 \left|\int a\,\mathrm df\wedge\dc\omega\wedge T\right|
 \leq C\left(\int a\,\mathrm df\wedge\dc f\wedge\omega\wedge T\right)^{1/2}
          \left(\int a\,\omega^2\wedge T\right)^{1/2}.
\]
The corresponding inequality with $\dc f\wedge\mathrm d\omega$
holds as well. For a mixed product of $q$ cone forms in place of
$S^q$, it suffices to sum these bounds over the finitely many
subset sums of its factors.
\end{lemma}
\begin{proof}
Diagonalize $S$ in an $\omega$-unitary frame. Every diagonal
coefficient of $S^q\omega^{n-q-2}$ is a positive multiple of
$S_{q;ij}(\lambda)$ for two omitted coordinates. This is nonnegative
by deleted-coordinate positivity, because $q+2\leq m$.
The form is a nonnegative sum of simple positive coordinate forms.

For one such simple form, only the remaining complex two-plane
contributes to a complementary wedge. The coefficients of
$\dc\omega$ are uniformly bounded in unitary frames on the fixed
compact set. Its wedge with $\mathrm df$ is consequently bounded
by a constant times the gradient norm on that plane. The wedges
with $\omega^2$ and with $\mathrm df\wedge\dc f\wedge\omega$
give, up to fixed positive constants, one and the square of that
norm. Pointwise Cauchy--Schwarz over the simple summands, followed
by integral Cauchy--Schwarz with weight $a$, proves the estimate.
Conjugation gives the other estimate.

For $q\geq1$, use the exact polarization identity
\[
 q!\,\alpha_1\wedge\cdots\wedge\alpha_q
   =\sum_{J\subset\{1,\ldots,q\}}
           (-1)^{q-|J|}\left(\sum_{i\in J}\alpha_i\right)^q.
\]
Every nonempty subset sum remains in the cone. Applying the single
power estimate to each term and taking absolute values proves the
mixed assertion. At $q=0$ use the single term $S^0=1$.
No Cauchy--Schwarz assertion for an arbitrary weakly positive
$(n-2,n-2)$-form is needed.
\end{proof}

\begin{lemma}\label{draft621:boundary-localization}
Let $f_j$ be uniformly bounded smooth local $m$-subharmonic functions
with $f_j\downarrow f$, and fix a smaller concentric ball $V\Subset U$.
There are uniformly bounded smooth local functions $v_j\downarrow v$
which equal $f_j$ on a neighborhood of $\overline V$ and equal one
fixed strictly plurisubharmonic $\psi$ on one fixed boundary collar.
Moreover $v=f$ near $\overline V$ and $v=\psi$ on that collar.
\end{lemma}
\begin{proof}
Write $U=\{|z|<R\}$ and $|f_j|\leq M$. Take
$\psi=A(|z|^2-R^2)+b$, with $b>M+2$ and then $A$ so large that
$\psi<-M-2$ on a neighborhood of $\overline V$.
It exceeds $M+2$ on a fixed sufficiently thin boundary collar.

Let $\zeta$ be a nonnegative even smooth probability density
supported in $[-1,1]$, and put
\[
 \chi_\epsilon(t)=\int|t+\epsilon s|\zeta(s)\,\mathrm ds,
 \qquad
 M_\epsilon(a,b)=\tfrac12(a+b+\chi_\epsilon(a-b)).
\]
This is smooth and convex, increases in each argument, and equals
$\max(a,b)$ when $|a-b|\geq\epsilon$. It is nondecreasing in
$\epsilon$: pairing $s$ and $-s$ shows that the scale derivative
of $\chi_\epsilon$ is nonnegative. Choose $0<\epsilon_j\downarrow0$
with $\epsilon_j<1$, and put $v_j=M_{\epsilon_j}(f_j,\psi)$.
Both monotonicities give $v_j\downarrow\max(f,\psi)=v$.
The fixed gaps give the two asserted regions of exact equality.

Differentiation writes $\ddc v_j$ as a convex combination of
$\ddc f_j$ and $\ddc\psi$, plus
$\tfrac12\chi_{\epsilon_j}''(f_j-\psi)
       \mathrm d(f_j-\psi)\wedge\dc(f_j-\psi)$.
The last form is nonnegative. Cone convexity and stability under
nonnegative addition prove local $m$-subharmonicity. Bounds for
$f_j,\psi$ and $\chi_\epsilon(t)-|t|\leq\epsilon$ give the common
bound. In particular all differences $v_k-v_j$ have support in
one fixed compact subset of $U$.
\end{proof}

\begin{proposition}\label{draft621:smooth-capacity-convergence}
For the localized sequence in the preceding lemma and
$0\leq s\leq m-1$,
\[
 c_s\{v_k-v>\delta\}\longrightarrow0
                 \qquad(\delta>0).
\]
Consequently a uniformly bounded smooth decreasing local sequence
has a quasicontinuous pointwise limit on every compact subset of
$U$, for $c_{m-1}^*$ as well as $c_{m-1}$.
\end{proposition}
\begin{proof}
The assertion for $s=0$ follows by bounded convergence for volume.
Induct on $s$, assuming it for $s-1$, and put $q=s-1$ and
$\theta=\omega^{n-s}$. For $j\geq k$ let $h=v_k-v_j\geq0$.
The support of $h$ lies in one fixed compact set $K_0\Subset U$.
Smooth capacity domination and the induction hypothesis give,
uniformly over all bounded local smooth mixed potentials,
\begin{equation}\label{eq:draft621:lower-weight}
 \int h\,\ddc g_1\cdots\ddc g_q\omega^{n-q}
                          \longrightarrow0\quad(k\to\infty).
\end{equation}
To see this, split at $h=\delta$. On $\{h\leq\delta\}$ the mass
bound gives $C\delta$. On the remaining set use $h\leq v_k-v$,
the capacity bound and the induction hypothesis. Let $k$ tend to
infinity before $\delta\downarrow0$. For $q=0$ this is the volume
estimate. Constants are uniform in $j\geq k$ and all potentials
with the fixed common bound.

Consider the real smooth energy
\[
 E_s(w)=\int_U w(\ddc w)^s\theta,
 \qquad w_t=v_j+t h,\quad0\leq t\leq1.
\]
The energies $E_s(v_j)$ are bounded: on the fixed collar the
integrand is fixed, and the remaining compact set has a uniform
smooth mass bound. Since $h$ vanishes near the boundary,
integration by parts and closedness of $\ddc w_t$ give the exact
identity
\begin{equation}\label{eq:draft621:energy-variation}
 \begin{split}
 \frac{\mathrm d}{\mathrm dt}E_s(w_t)
  &=(s+1)\int h(\ddc w_t)^s\theta\\
  &\quad+s\int h(\ddc w_t)^q
       \bigl(\mathrm dw_t\wedge\dc\theta
           -\dc w_t\wedge\mathrm d\theta+w_t\ddc\theta\bigr).
 \end{split}
\end{equation}
Here
\[
 \mathrm d\theta=(n-s)\mathrm d\omega\wedge\omega^{n-s-1},
 \quad
 \ddc\theta=(n-s)\ddc\omega\wedge\omega^{n-s-1}
 +(n-s)(n-s-1)\mathrm d\omega\wedge\dc\omega
                                           \wedge\omega^{n-s-2},
\]
and the analogous $\dc$ formula holds. A zero scalar coefficient
means that its term is absent. The two first-order terms in
\eqref{eq:draft621:energy-variation} tend uniformly to zero by
\cref{draft621:smooth-torsion}, the uniform local gradient bound,
and \eqref{eq:draft621:lower-weight}. The last term has the same
property by the absolute-form estimate at degree $q<m$.
More explicitly, the first torsion integral is bounded by
\[
 C\left(\int h\,\mathrm dw_t\wedge\dc w_t
                     \wedge(\ddc w_t)^q\omega^{n-q-1}\right)^{1/2}
   \left(\int h(\ddc w_t)^q\omega^{n-q}\right)^{1/2}.
\]
The first factor is bounded since $h$ has a common bound and
\cref{draft621:local-smooth-bounds} applies; the second tends to
zero by \eqref{eq:draft621:lower-weight}. Conjugation handles the
other first-order term.
Thus, with $e_k\downarrow0$ after taking a tail supremum,
\[
 E_s(v_k)-E_s(v_j)=(s+1)\Delta_{kj}+r_{kj},
 \quad
 \Delta_{kj}=\int_0^1\int h\Lambda_s(w_t)\,\mathrm dt\geq0,
 \quad |r_{kj}|\leq e_k.
\]
It follows that $E_s(v_j)\leq E_s(v_k)+e_k$ for $j\geq k$.
Taking a subsequence $k$ approaching the liminf shows that the
bounded energy sequence converges. Hence $\Delta_{kj}\to0$
uniformly for $j\geq k$. Positivity of the binomial expansion
$\Lambda_s(w_t)\geq(1-t)^s \Lambda_s(v_j)$ then gives
\begin{equation}\label{eq:draft621:self-weight}
 \int h\Lambda_s(v_j)\longrightarrow0
                         \quad\text{uniformly for }j\geq k.
\end{equation}

It remains to replace each Hessian factor by a capacity competitor.
For any smooth $-1\leq a\leq0$ in the local cone, put
\[
 I_\ell=\int h(\ddc v_j)^\ell(\ddc a)^{s-\ell}\theta,
 \qquad
 \eta_\ell=(\ddc v_j)^\ell(\ddc a)^{s-\ell-1},
                \quad0\leq\ell<s.
\]
First-order Stokes, with every metric derivative retained, gives
\[
 I_\ell=-\int\mathrm dh\wedge\dc a\wedge\eta_\ell\theta
          +\int h\,\dc a\wedge\eta_\ell\mathrm d\theta.
\]
The second term is $o(1)$ uniformly in $a$ and $j\geq k$.
Use finite polarization in \cref{draft621:smooth-torsion}; each
subset sum has a uniform local gradient bound and its lower trace
is controlled by \eqref{eq:draft621:lower-weight}.

The form $\eta_\ell\theta$ is positive of bidegree $(n-1,n-1)$:
a complementary nonnegative rank-one form completes a polarized
product of $s\leq m$ cone forms. Its associated Hermitian
gradient pairing is nonnegative, so Cauchy--Schwarz and the uniform
gradient bound for $a$ give
\[
 I_\ell\leq C F_\ell(h)^{1/2}+o(1),
 \qquad F_\ell(h)=\int\mathrm dh\wedge\dc h\wedge\eta_\ell\theta.
\]
Another exact first-order Stokes identity gives
\[
 F_\ell(h)= I_{\ell+1}
       -\int h\ddc v_k\wedge\eta_\ell\theta
       +\int h\,\dc h\wedge\eta_\ell\mathrm d\theta
                    \leq I_{\ell+1}+o(1).
\]
The dropped term is nonnegative. For the last term the same
polarized torsion estimate applies: its first gradient factor is
uniformly bounded because $h=v_k-v_j$, and its second factor
tends to zero by \eqref{eq:draft621:lower-weight}; here
$h^2\leq C h$ and the gradient energy of the difference is at
most twice the sum of those of $v_k$ and $v_j$ for every positive
single-power gradient pairing.
Starting with \eqref{eq:draft621:self-weight} at $I_s$, finite
backward induction gives $I_0\to0$ uniformly in $a$ and $j\geq k$.

Chebyshev's inequality now yields
$c_s\{v_k-v_j>\delta\}\leq\delta^{-1}\sup_a I_0$.
These sets increase to $\{v_k-v>\delta\}$ as $j\to\infty$.
Continuity of $c_s$ under increasing unions proves the induction.

Choose $k_N$ so that the open sets
$G_N=\{v_{k_N}-v>2^{-N}\}$ satisfy
$c_{m-1}(G_N)<\epsilon2^{-N}$, $N\geq1$.
They are open because $v$ is upper semicontinuous and $v_{k_N}$
is smooth. On $U\setminus\bigcup_NG_N$ these smooth functions
converge uniformly to the pointwise limit $v$. Hence $v$ is
continuous there. The union has capacity less than $\epsilon$.
Boundary localization transfers this conclusion to $f$ on each
smaller compact set. Open sets have the same outer capacity, which
also proves the assertion for $c_{m-1}^*$.
\end{proof}

\begin{lemma}\label{draft621:rough-coefficient-capacity}
Every bounded admissible potential on $X$ is locally quasicontinuous
for the preceding smooth capacity $c_{m-1}^*$. For $0\leq p<m$,
each coefficient measure $\nu$ of the book's ordered $D_p$ satisfies
\[
 |\nu|(E)\leq C c_{m-1}^*(E),\qquad E\subset K\Subset U,
\]
uniformly when its global potentials have one common bound.
The assertion at $p=0$ means the volume coefficient of $D_0=1$.
\end{lemma}
\begin{proof}
The global decreasing approximation and
\cref{hessian:local-unshifted-correction} provide bounded smooth
local approximations to $u+\rho$. The preceding proposition makes
its canonical pointwise limit quasicontinuous. Subtracting the
smooth $\rho$ preserves quasicontinuity.

For smooth approximations write
$\ddc u_i=\ddc(u_i+\rho)-\ddc\rho$ and expand the product.
Each term has $q\leq p<m$ local cone factors and fixed smooth
complementary factors. The absolute-form estimate and
\eqref{eq:draft621:smooth-capacity-bound}, with $s=m-1$, bound
each coefficient variation by $C c_{m-1}$ on Borel sets.
The constants depend only on the common bound and local data.

Take the iterated limits of
\cref{hessian:bounded-recursive-currents} on a slightly larger
compact neighborhood of $K$. Lower semicontinuity of signed
measure variation on open sets preserves these estimates at each
limit. For any open $G\supset E$, intersect with a fixed open
neighborhood of $K$ having compact closure in $U$; this gives
$|\nu|(E)\leq Cc_{m-1}(G)$. Taking the infimum proves the outer
estimate. No simultaneous rough product or rough diagonal
polynomial expansion is used.
\end{proof}

\begin{lemma}\label{draft621:weighted-measure-convergence}
Suppose signed Radon measures $\nu_j\to\nu$ weakly have locally
uniformly bounded variations, each dominated on compact sets by
$C c_{m-1}^*$. If bounded Borel functions $a_j\downarrow a$ have
one common bound and every member of this countable family is
locally quasicontinuous for that capacity, then
$a_j\nu_j\to a\nu$ weakly.
\end{lemma}
\begin{proof}
Fix $\phi\in C_c(U)$ and a compact $K\Subset U$ containing its
support. Summably allocate exceptional capacities less than
$\epsilon$ to the countable family. Their open union $G$ has
$c_{m-1}^*(G)<\epsilon$, and every function is continuous on
the compact set $F=K\setminus G$. Dini's theorem gives
$\|a_j-a\|_F\to0$. Extend $a|_F$ to a bounded continuous $g$ on
$U$ with the same bound $M$; if $F$ is empty take $g=0$.
Then
\[
 \begin{split}
 \left|\int\phi(a_j\,\mathrm d\nu_j-a\,\mathrm d\nu)\right|
 &\leq\|\phi\|_\infty\|a_j-a\|_F|\nu_j|(K)
          +\left|\int\phi g\,\mathrm d(\nu_j-\nu)\right|\\
 &\quad+2M\|\phi\|_\infty
                 \bigl(|\nu_j|(K\cap G)+|\nu|(K\cap G)\bigr).
 \end{split}
\]
For fixed $G$ the first two terms tend to zero, and the last is
at most $4MC\|\phi\|_\infty\epsilon$. Let $\epsilon\downarrow0$.
The weak-measure passage uses a continuous extension on the whole
measure, rather than the discontinuous indicator of $F$.
\end{proof}

\begin{theorem}\label{draft621:joint-decreasing-products}
For uniformly bounded admissible potentials $u_i^j\downarrow u_i$
on $X$, the currents $D_p(u_1^j,\ldots,u_p^j)$ converge to
$D_p(u_1,\ldots,u_p)$ in distributions for $p\leq m$, and weakly
with measure coefficients for $p<m$. These currents are symmetric
in their potentials. For $p<m$ the coefficient convergence also
holds after multiplication by a uniformly bounded decreasing
admissible weight. The positive shifted currents $Q_p$ converge
weakly for $p\leq m$; $Q_m$ converges as a positive Radon measure.
No unwedged order-zero assertion for $D_m$, or weighted convergence
at degree $m$, is part of the conclusion.
\end{theorem}
\begin{proof}
Work on a compact coordinate test set. Degree zero is constant.
Suppose the result holds for $D_{p-1}$, which has measure
coefficients since $p-1<m$. The preceding two lemmas, applied
coefficientwise with weight $u_p^j$, give convergence of
$u_p^jD_{p-1}(u_1^j,\ldots,u_{p-1}^j)$. Applying $\ddc$ proves
distributional convergence at degree $p$. For $p<m$, the uniform
coefficient bounds in \cref{hessian:bounded-recursive-currents}
upgrade this to weak convergence of measures, by smooth
approximation of compactly supported continuous tests. The same
weighted lemma gives the additional weighted conclusion.

For each fixed list choose simultaneous smooth decreasing global
approximations. Their smooth raw Hessian wedges are symmetric.
The convergence just proved identifies their limit with $D_p$
in any ordering, proving symmetry without an exchange of
previously unproved iterated limits.

The subset expansion for $R_p$ uses fixed smooth powers of
$\omega$, so its distributional convergence and symmetry follow
from those of $D_I$. Its wedge $Q_p$ is positive with uniform
local masses by \cref{hessian:bounded-shifted-positive}.
Distributional convergence and these mass bounds imply weak
convergence as an order-zero positive current, including the
Radon measure conclusion at $p=m$. A partition of unity on a
finite coordinate cover proves the global assertions.
\end{proof}
